\scanpage{52}
% archive foliation: 51
configuration \deleted{\unsure{icosaèdre}} \inserted{octaédrale} associée est bien celle de départ\ldots

\uline{Théorème}\quad Soit $X$ un fibré en droites projectives sur un schéma $S$, de car. résiduelles $\ne2$. Considérons le \deleted{\unsure{gr.}} schéma en groupes $G=\underline{\operatorname{Aut}}_S(X)$ (\uline{NB} la donnée de $X/S$ équivaut à celle de la ``forme'' $G$ de $(\operatorname{PGL})_S$ sur $S$), et les ensembles suivants :

 a) L'ens $A(X)$ des sous-schémas en groupes $V$ de $G$, loc isom. à $(\mathbb D_4)_S$.

 b) L'ens \deleted{\unsure{de}} $B(X)$ des sous-schémas en groupes \deleted{\unsure{$G$}} $\Gamma$ de $G$, loc. isom. à $(\mathfrak S_4)_S$.

 c) L'ens des configurations icosaédrales régulières dans $X$, i.e. des diviseurs étales $\Delta$ de $X/S$, de degré $6$, munis d'une opposition $\underline a_\Delta:\Delta-\Delta$ qui \deleted{\unsure{définit} des pts} fixes sur chaque fibre géom., et tel que $\underline{\operatorname{Aut}}^+_{S\mathrm{oct}}(\Delta)$ opère sur $X$ de façon compatible avec son opération sur $\Delta$.

On a alors un système transitif de bijections canoniques entre ces trois ensembles, \unsure{défini} ainsi !
